\subsection{A3: an explicit very ample surface with negative first cohomology}

\begin{theorem}[A3; Proved]\label{a:very-ample-example}
There is a smooth geometrically integral projective surface $Y$ over an
algebraically closed field of characteristic two, and a very ample
invertible sheaf $L$ on $Y$, for which
\[
 h^0(Y,L)=562,\qquad h^1(Y,L^{-1})=1,\qquad h^1(Y,\OO_Y)=565.
\]
Thus the least dimension in the range requested in A3, namely
$\dim X\geq2$, is two. For the geometric generic hyperplane of this
embedding the kernel on $\mathbf{Pic}^{\tau}$ is nontrivial and has
one-dimensional tangent space. This conclusion concerns the group
scheme kernel; it does not give a counterexample on geometric points
to A1 or A2.
\end{theorem}

We give all the defining data, including a finite algorithm specifying
every coefficient of the sections used in the construction. The final
ground field is an explicitly specified algebraic closure of a purely
transcendental extension of $\mathbb F_2$. No numerical search for smooth
branch divisors is substituted for the proof of smoothness.

\paragraph{1. The curve and the extension.}
Let $C$ be the smooth projective model over $\mathbb F_2$ of
\[
 y^2+y=x^{15}.
\]
The Artin--Schreier polynomial is irreducible over
$\overline{\mathbb F}_2(x)$: a rational function whose Artin--Schreier
difference has a pole has an even pole order there, whereas $x^{15}$
has an odd pole order at infinity. There is one point $\infty$ over
infinity, with $v_\infty(x)=-2$ and $v_\infty(y)=-15$.
Write $P_0=(0,0)$, $P_1=(0,1)$, and put $t=y/x^8$. Then
\[
 \operatorname{div}(t)=7P_0+\infty-8P_1,\qquad
 x+t^{-2}=t^{-1}x^{-7},\qquad dt=x^6\,dx.
\]
In particular $t$ is a uniformizer at $\infty$, and
$v_\infty(x+t^{-2})=13$. The affine curve is smooth and $dx$ is
nowhere zero there, since the derivative of its equation with respect
to $y$ is one. The displayed formula for $dt$ therefore gives
\[
 \operatorname{div}(dx)=12\infty,\qquad g(C)=7.
\]
These formulas also verify the extension and smoothness calculations
below directly, without an appeal to an unspecified Tango curve.

Set $N=\OO_C(2\infty)$ and $D=6\infty$, so that
$\OO_C(D)=N^3$. Cover $C$ by
$U=C\setminus\{\infty\}$ and $V=C\setminus\{P_0,P_1\}$.
Define a rank-two bundle $E$ by frames $(e_0,e_1)$ on $U$ and
$(v_0,v_1)$ on $V$, with transition
\begin{equation}\label{a:explicit-transition}
 v_0=e_0,\qquad v_1=t^{-6}(e_1-t^{-1}e_0).
\end{equation}
It fits into $0\to\OO_C\to E\to N^3\to0$.
The extension is nonsplit. Indeed a splitting would express $t^{-1}$
as a function regular on $U$ plus a function divisible by $t^6$ and
regular at $\infty$. The function on $U$ would have its sole pole,
of order one, at $\infty$. This is impossible: the ring of $U$ is
$\mathbb F_2[x,y]/(y^2+y-x^{15})$, and its possible pole orders are
$2a$ and $15+2a$, $a\geq0$.

The Frobenius pullback splits. On $U$ its second summand is generated
by $e_1^2+xe_0^2$; on $V$ the corresponding expression is
\[
 v_1^2+t^{-12}(x+t^{-2})v_0^2.
\]
The coefficient in this expression is regular on $V$ and has order
one at $\infty$. This gives a splitting of the quotient $N^6$ and
also a map $N^6\to\operatorname{Sym}^2E$.

\paragraph{2. The smooth cyclic-cover surface.}
Use the quotient convention $P=\mathbf P_C(E)$, so that
$\pi_*\OO_P(j)=\operatorname{Sym}^jE$ for $j\geq0$.
Let $E_\infty$ be the section $e_0=0$; its divisor sheaf is
$\OO_P(1)$ and $E_\infty^2=6$. The map just constructed defines
a divisor
\[
 C''\in\big|\OO_P(2)\otimes\pi^*N^{-6}\big|.
\]
On $U$, in the affine chart $e_0\ne0$ with $w=e_1/e_0$, its
equation is $w^2+x=0$. It is disjoint from $E_\infty$.
It is smooth on this chart since $dx$ never vanishes on $U$.
Over $\infty$ its equation is
$w_V^2+t^{-12}(x+t^{-2})=0$; the second term has order one,
so the same smoothness assertion holds there. Thus $C''$ and
$E_\infty$ are disjoint smooth divisors.

Put $M=\OO_P(-1)\otimes\pi^*N^2$. Then
$M^{-3}=\OO_P(E_\infty+C'')$. Let $\sigma$ be the product of
the two defining sections, and form the cyclic cover
\[
 \psi:S=\operatorname{Spec}_P
       (\OO_P\oplus M\oplus M^2)\longrightarrow P,
 \qquad M^3\xrightarrow{\sigma}\OO_P.
\]
This is finite flat of degree three. Locally at the branch locus it
has equation $z^3=u$ with $u$ a regular parameter. Since three is
invertible in characteristic two and the branch components are smooth
and disjoint, $S$ is smooth. It is projective because it is finite over
$P$. Moreover $H^0(S,\OO_S)=\mathbb F_2$: the two nonconstant
summands have negative degree on the fibres of $\pi$. This remains
true after extension of the ground field, and smoothness implies
geometric integrality. Write $f=\pi\psi$ and let $T$ be the
ramification section over $E_\infty$. Then
\[
 \psi^*E_\infty=3T,\qquad T\simeq C,\qquad T^2=2.
\]
On the dense affine chart used above all equations are simply
\begin{equation}\label{a:surface-affine}
 y^2+y=x^{15},\qquad z^3=w^2+x.
\end{equation}

\paragraph{3. An ample generated bundle with the required cohomology.}
Define
\[
 Q=\OO_S(2T)\otimes f^*N.
\]
The bundle $E$ is nef: after pullback to any smooth projective curve,
a quotient line bundle receives either a nonzero map from $\OO$ or
factors through the positive-degree quotient $N^3$; in either case
its degree is nonnegative. Thus $\OO_P(1)$ is nef. The line bundle
$\OO_P(2)\otimes\pi^*N^3$ has positive intersection with every
curve and positive self-intersection: on a fibre its degree is two,
and on a horizontal curve its base factor has positive degree.
The Nakai criterion makes it ample. Since
\[
 Q^3=\psi^*(\OO_P(2)\otimes\pi^*N^3),
\]
$Q$ is ample as well.

For an integer $a$, the local equation $z^3=u$ at $T$ gives
\begin{equation}\label{a:cyclic-floor}
 \psi_*\OO_S(aT)=
 \bigoplus_{i=0}^2 M^i\bigl(\lfloor(a+i)/3\rfloor E_\infty\bigr).
\end{equation}
Indeed the $i$th eigensummand consists locally of
$z^i u^{-\lfloor(a+i)/3\rfloor}$ times regular functions.
For $a=2$, this formula and projection to $C$ give
\[
 f_*Q=N\oplus N^3.
\]
Both summands are generated by global sections: their bases are
$(1,x)$ and $(1,x,x^2,x^3)$, respectively. The first summand generates
$Q$ away from $T$. The second has valuation $1-3+2=0$ along $T$
and generates there. Hence $Q$ is globally generated and
\[
 h^0(S,Q)=6;\qquad
 1,x,z,zx,zx^2,zx^3
\]
are its sections in the affine trivialization of
\eqref{a:surface-affine}.

For $a=-2$, formula~\eqref{a:cyclic-floor} instead gives
\[
 \psi_*Q^{-1}=
 \bigl(\OO_P(-1)\otimes\pi^*N^{-1}\bigr)
 \oplus\bigl(\OO_P(-2)\otimes\pi^*N\bigr)
 \oplus\bigl(\OO_P(-2)\otimes\pi^*N^3\bigr).
\]
Relative duality, with the convention stated above, says
$R^1\pi_*\OO_P(-2)=(\det E)^{-1}=N^{-3}$.
All zeroth direct images in the preceding display vanish, so Leray
gives
\begin{equation}\label{a:Q-cohomology}
 H^1(S,Q^{-1})=H^0(C,N^{-2}\oplus\OO_C)=\mathbb F_2.
\end{equation}
Proposition~\ref{a:raynaud-all}, applied with $a=2n$, $b=n$,
$p=2$ and $\ell=3$, also gives
\[
 H^1(S,Q^{-n})=0\quad(n\geq2).
\]
For later dimension checks, $R^1f_*\OO_S=N$, whence
$h^1(S,\OO_S)=7+h^0(C,N)=9$. Intersection and duality give
$Q^2=16$, $K_S=f^*N^5$, $K_SQ=20$, and $\chi(\OO_S)=-2$.
The canonical formula follows as well from
$K_P=\OO_P(-2)\otimes\pi^*(K_C\otimes\det E)$ and
$K_S=\psi^*(K_P\otimes M^{-2})$.

\paragraph{4. A specified very ample power and its complete basis.}
Set $B=\OO_P(6)\otimes\pi^*N^9$. Then
\[
 Q^9=\psi^*B,\qquad
 \psi_*Q^9=B\oplus BM\oplus BM^2
   =\bigoplus_{i=0}^2\OO_P(6-i)\otimes\pi^*N^{9+2i}.
\]
The bundle $\operatorname{Sym}^a E\otimes N^b$ has a filtration
with quotients $N^{b+3j}$, $0\leq j\leq a$.
For $(a,b)=(6,9),(5,11),(4,13)$ all these quotients, even after
subtracting any effective divisor of degree two on $C$, have degree
at least sixteen, greater than $2g(C)-2=12$.
Thus their first cohomology vanishes, as does that of the filtered
bundle. Restriction of sections is consequently surjective over every
base subscheme of length at most two.

It follows that $B$ is very ample on $P$. To see every infinitesimal
case explicitly, a length-two subscheme of $P$ either lies over a
length-two subscheme of $C$, or lies in one fibre. In the first case
use the just-proved surjectivity and relative very ampleness of
$\OO_P(6)$; in the second, use $\OO_{\mathbf P^1}(6)$.
The same argument with length one proves that $BM$ is globally
generated. These facts show that $\psi^*B$ is very ample on $S$:
the $B$ sections recover the embedded base $P$, while, on an open
where a section of $BM$ is nonzero, its corresponding section of
$\psi^*B$ is a unit multiple of the cyclic generator $z$.
They generate the local algebra of the degree-three cover and hence
give a closed immersion. Equivalently they separate length-two
subschemes whose image is a point as well as those separated on $P$.

Riemann--Roch and the same vanishing give the exact dimensions
\[
 h^0(P,B)=210,\quad h^0(P,BM)=186,\quad
 h^0(P,BM^2)=160,\quad h^0(S,Q^9)=556.
\]
For complete explicit defining coefficients, the ordered basis
$v_1,\ldots,v_{556}$ is the basis over $\mathbb F_2$ in
\texttt{code/A3\_Q9\_basis.json}, generated by
\texttt{code/certify\_A3.py}. Here is its finite mathematical
specification. For $i=0,1,2$ put $a=6-i$, $b=9+2i$, and order the
monomials $w^jx^ry^sz^i$ lexicographically by $(j,s,r)$, where
\[
 0\leq j\leq a,\quad s\in\{0,1\},\quad r\geq0,
 \qquad 2r+15s\leq2b+7a-j.
\]
Use the expansion $x=t^{-2}u$, $y=t^{-15}u^8$, where
$u=1+t^{15}u^{-7}$. Impose that, for each $0\leq h\leq a$,
every negative power of $t$ in
\[
 \sum_{j\geq h,r,s}c_{jrs}{j\choose h}
       t^{2b+7h-j-2r-15s}u^{r+8s}
\]
has coefficient zero. This is precisely the regularity condition
after $w_U=t^6w_V+t^{-1}$ and $z_U=t^4z_V$.
Order the equations by their first occurrence when the ordered
monomials are processed, and take the reduced free-variable kernel
basis with free variables in increasing order, as specified by the
included elimination routine. Concatenate the three bases in
increasing $i$. All sections are regular on the affine curve, on
the chart at $\infty$, and along the section $T$ by their indicated
projective powers. Independence and the Riemann--Roch dimensions prove
that they form the complete basis.

The exact calculation has respectively $357,291,230$ variables,
ranks $147,105,70$, and nullities $210,186,160$. Only four, four,
and three terms beyond the constant in the $t^{15}$ expansion can
contribute to the negative powers, so the calculation is finite and
exact. Its output and time are recorded in
\texttt{code/certify\_A3\_output.json}; the basis SHA256 is recorded
there too. These checks verify coefficient and transition equations,
not a sampled collection of points on the surface.

\paragraph{5. The cover making the first power very ample.}
Let
\[
 k=\overline{\mathbb F_2(\lambda_{ij}\mid1\leq i,j\leq556)},
 \qquad s_i=\sum_{j=1}^{556}\lambda_{ij}v_j.
\]
All varieties and bundles above are now extended to $k$. The matrix
$\lambda$ is invertible, so the $s_i$ are a basis of $H^0(S,Q^9)$.
Their zero divisors are smooth, pairwise transverse, and no three
meet. Here is the exact genericity argument. The smooth-divisor
Bertini theorem applies because $Q^9$ is very ample (Stacks,
Lemma~33.47.3, opened)~\cite{StacksBertini}. After selecting any number of such divisors,
apply the same theorem to their smooth components to make the next
divisor transverse, and avoid their finitely many pairwise
intersection points. Also avoid the linear span of the already
chosen sections. At each stage these are finitely many nonempty
open conditions. Their intersection is nonempty over an infinite
field. They define a nonempty open set of ordered bases; algebraic
independence of all $\lambda_{ij}$ places our basis in this set.

For each $s_i$ take the ninth-root algebra
$\bigoplus_{r=0}^{8}Q^{-r}$ with multiplication determined by
$s_i:Q^{-9}\to\OO_S$, and set
\begin{equation}\label{a:root-algebra}
 \mathcal A=\bigotimes_{i=1}^{556}
       \left(\bigoplus_{r=0}^{8}Q^{-r}\right),\qquad
 \rho:Y=\operatorname{Spec}_S\mathcal A\longrightarrow S.
\end{equation}
This describes the whole finite cover, not a choice of a connected
component. It has degree $9^{556}$. There are tautological sections
$t_i$ of $L=\rho^*Q$ with $t_i^9=\rho^*s_i$.

The surface $Y$ is smooth. At any point of $S$ at most two branch
divisors pass, with independent regular parameters as equations.
After an étale extension extracting ninth roots of units, the cover
is obtained by replacing each such parameter by its ninth power.
The other equations define étale extensions, since nine is
invertible. The resulting local rings are regular of dimension two.
Projectivity follows from finiteness. Moreover
\begin{equation}\label{a:root-directsum}
 \rho_*\OO_Y=
 \bigoplus_{0\leq r_1,\ldots,r_{556}\leq8}
                   Q^{-(r_1+\cdots+r_{556})}.
\end{equation}
Ampleness of $Q$ makes all negative-power global sections zero, so
$H^0(Y,\OO_Y)=k$. Thus $Y$ is connected, and its smoothness over the
algebraically closed field makes it integral.

We verify very ampleness of $L$ without conflating this cover with
the inverse image under a coordinate power map. Use the six pulled
back sections of $Q$ and the $556$ sections $t_i$. They generate $L$
because the old six already do. Let $u$ be one of the old sections.
The open $S_u=S\setminus V(u)$ is affine: it is the intersection of
the $Q^9$ embedding with the affine hyperplane chart $u^9\ne0$.
Since the $s_i$ form a complete basis, its coordinate algebra is
generated over $k$ by $s_i/u^9$. Over this open put $z_i=t_i/u$.
Then
\[
 \OO(Y_u)=\OO(S_u)[z_1,\ldots,z_{556}]
                   /(z_i^9-s_i/u^9)_{i=1}^{556}.
\]
This algebra is generated over $k$ by the $z_i$ alone, because their
ninth powers include all the displayed generators of $\OO(S_u)$.
Consequently the projective morphism defined by the old sections and
the $t_i$ is a closed immersion on each target chart $u\ne0$.
Those charts cover $Y$. The morphism is therefore an immersion,
and properness makes it a closed immersion globally. Hence $L$ is
very ample. Retaining the old sections in this argument also
separates simultaneous multiplication of all $t_i$ by a ninth root
of unity; no torsion ambiguity in the identity $L=\rho^*Q$ remains.

Finally tensor~\eqref{a:root-directsum} by $Q^{-1}$ and use
\eqref{a:Q-cohomology} and the higher negative-power vanishing.
Only the multi-index $(0,\ldots,0)$ contributes to $H^1$, giving
$h^1(Y,L^{-1})=1$. The same decomposition gives
$h^0(Y,L)=6+556=562$ and $h^1(Y,\OO_Y)=9+556=565$.
This completes all geometric and cohomological assertions in
Theorem~\ref{a:very-ample-example}.

\paragraph{Verification and attribution.}
Three other AI agents independently checked the root-cover argument,
the smoothness and connectedness assertions, the cohomology summands,
and the very-ampleness chart proof; their audit notes are retained.
The finite certificate is independent of those prose checks. The
construction uses the classical Raynaud mechanism, but the equations,
specified polarization and simultaneous-root proof above are supplied
in full here. Searches for a published theorem with this exact
very-ample negative-$H^1$ conclusion did not locate one during this
run. No priority claim is made.
